% ============================================================
%  Paper 1 — Aroha Labs
%  "Asymmetric Temporal Decay in Bayesian Reputation Systems
%   for Open Multi-Agent AI Networks"
%
%  Target venues: AAMAS 2027 / IJCAI 2027 / NeurIPS 2026 Workshop
%  arXiv categories: cs.MA, cs.AI, cs.LG
% ============================================================

\documentclass[conference]{IEEEtran}
\IEEEoverridecommandlockouts

\usepackage{amsmath,amssymb,amsthm}
\usepackage{mathtools}
\usepackage{algorithm}
\usepackage{algpseudocode}
\usepackage{graphicx}
\usepackage{booktabs}
\usepackage{hyperref}
\usepackage{xcolor}
\usepackage{cite}
\usepackage{bm}

\newtheorem{theorem}{Theorem}
\newtheorem{lemma}[theorem]{Lemma}
\newtheorem{proposition}[theorem]{Proposition}
\newtheorem{corollary}[theorem]{Corollary}
\newtheorem{definition}{Definition}
\newtheorem{remark}{Remark}

\DeclareMathOperator{\E}{\mathbb{E}}
\DeclareMathOperator{\Var}{\mathrm{Var}}
\DeclareMathOperator{\Beta}{\mathrm{Beta}}
\newcommand{\R}{\mathbb{R}}

\begin{document}

\title{Asymmetric Temporal Decay in Bayesian Reputation Systems\\
       for Open Multi-Agent AI Networks}

\author{
  \IEEEauthorblockN{Aroha Labs Research}
  \IEEEauthorblockA{%
    \textit{Aroha Protocol}\\
    aroha-labs.com
  }
}

\maketitle

% ─────────────────────────────────────────────────────────────
\begin{abstract}
Open multi-agent AI networks—environments in which autonomous
agents published by different organisations interoperate and
transact—require a principled mechanism for trust inference.
Existing reputation systems designed for peer-to-peer file
sharing or e-commerce fail to account for two properties
unique to AI agent networks: (i) observed agent behaviour
changes over time as models are updated or replaced, and
(ii) the cost of a false-positive trust assignment is
asymmetric with the cost of a false-negative, because a
malicious agent can cause real financial harm.
We present \textsc{ArohaRep}, a multi-signal Bayesian
reputation engine built on the Beta distribution that
introduces \emph{asymmetric temporal decay}: success evidence
decays with a half-life of $\tau_s = 180$ days while failure
evidence decays with a half-life of $\tau_f = 60$ days.
We formalise an eleven-signal observation model weighted by
signal reliability and combine Thompson Sampling with a
peer-endorsement gate to produce unbiased discovery rankings.
We prove that asymmetric decay strictly dominates symmetric
decay in expected regret under a model of agents with
time-varying quality, and provide closed-form bounds on
time-to-detection of quality degradation. Simulation
experiments over synthetic agent populations (1{,}000 runs,
code included) quantify the operator's trade-off: relative to
a symmetric-decay baseline, ArohaRep roughly halves
false-positive demotion of genuinely recovered agents
($18.6\%$ vs.\ $35.0\%$ of post-recovery agent-days) at the
cost of slower detection of degrading agents ($27.1$ vs.\
$15.9$ days). A sensitivity sweep shows the half-life ratio
$\rho = \tau_s/\tau_f$ is a direct dial on this trade-off,
with diminishing recovery gains beyond the deployed
$\rho = 3$. We argue this recovery-favouring operating point
is the correct one for Aroha specifically because financial
exposure is independently bounded by cryptographic spending
mandates, leaving discovery coverage as the reputation
system's binding cost.
\end{abstract}

\begin{IEEEkeywords}
multi-agent systems, reputation, Bayesian inference, Thompson
Sampling, temporal decay, trust, AI agents
\end{IEEEkeywords}

% ─────────────────────────────────────────────────────────────
\section{Introduction}
\label{sec:intro}

The emergence of autonomous AI agents capable of executing
multi-step tasks, calling external APIs, and handling
financial transactions has created a new class of distributed
system in which agents published by independent organisations
dynamically discover and invoke one another.
Unlike traditional microservice registries—where deployment
is controlled and services are assumed correct by construction
—open agent registries must cope with heterogeneous trust
levels, adversarial participants, and the fact that agent
quality changes over time as the underlying models are
updated, fine-tuned, or replaced.

The core question addressed by this paper is: \emph{given a
stream of noisy quality observations for an agent, how should
past observations be discounted over time to produce the most
accurate present estimate of quality?}

Classic solutions from peer-to-peer reputation literature
\cite{josang2002beta,kamvar2003eigentrust} apply either no
temporal decay or symmetric exponential decay, where all
historical evidence fades at a uniform rate. We identify two
fundamental problems with symmetric decay in the AI agent
setting:

\textbf{Problem 1 (Vulnerability to targeted degradation).}
With symmetric decay, a sustained campaign of low-quality
signal submissions can permanently damage a good agent's
reputation, because the failure evidence decays at the same
rate as the accumulated success evidence.

\textbf{Problem 2 (Slow recovery from transient failure).}
If an agent experiences a bad period due to a model rollback
or infrastructure issue, symmetric decay keeps the failure
evidence alive for as long as the pre-failure success
evidence. This penalises agents that have genuinely improved.

We propose \emph{asymmetric decay}: success pseudo-counts
decay slowly (half-life $\tau_s$) while failure pseudo-counts
decay quickly (half-life $\tau_f$), with $\tau_s > \tau_f$.
The intuition is that a long track record of good performance
constitutes stronger evidence of reliability than a recent
bad period, and should therefore be more resistant to decay.
Simultaneously, allowing failures to fade faster enables
genuine recovery by improved agents.

Our contributions are:
\begin{enumerate}
  \item A formal multi-signal observation model mapping eleven
    heterogeneous quality signals to Beta-distribution
    pseudo-count updates (\S\ref{sec:model}).
  \item The asymmetric decay operator $\mathcal{D}_{\tau_s,
    \tau_f}$ applied to Beta parameters, with a proof that
    it dominates symmetric decay under expected regret
    (\S\ref{sec:decay}).
  \item A peer-endorsement gate that prevents low-quality
    agents from influencing one another's reputation through
    colluding endorsements (\S\ref{sec:pps}).
  \item Thompson Sampling as the discovery ranking mechanism,
    with a regret bound for the agent-selection problem
    (\S\ref{sec:thompson}).
  \item Simulation results comparing ArohaRep to symmetric
    decay and no-decay baselines (\S\ref{sec:eval}).
\end{enumerate}

% ─────────────────────────────────────────────────────────────
\section{Related Work}
\label{sec:related}

\subsection{Beta Reputation Systems}

The foundational work by J{\o}sang and Ismail \cite{josang2002beta}
introduced the Beta distribution as a natural model for binary
reputation: a prior $\Beta(\alpha_0, \beta_0)$ is updated by
observed successes and failures, with the posterior mean
serving as a trust value. Their framework treats all
observations as equally weighted and does not model temporal
decay. We extend this framework with signal weighting and
asymmetric decay.

\subsection{Temporal Discount in Reputation}

Several extensions to the Beta model incorporate temporal
discount. Noorian and Ulieru \cite{ziegler2005propagation}
apply a sliding window that discards observations older than
a threshold. Yu and Singh \cite{yu2002peertrust} weight
recent interactions more heavily via a linear discount factor.
Both approaches use symmetric discount—all evidence ages at
the same rate—and neither considers the asymmetric cost
structure of success versus failure evidence. Our work is the
first to formally justify and quantify the benefit of
asymmetric decay rates for the specific cost structure of
AI agent quality assessment.

\subsection{Multi-Criteria Trust}

EigenTrust \cite{kamvar2003eigentrust} aggregates direct and
transitive trust via eigenvector computation. FIRE
\cite{sabater2005review} combines direct experience,
witnessed reputation, role-based knowledge, and third-party
certification. Neither handles temporal decay or the
asymmetric signal reliability of AI quality metrics. We
build a lighter-weight model appropriate for sub-100ms
latency requirements of agent discovery endpoints.

\subsection{Bandit Algorithms for Discovery}

Thompson Sampling \cite{thompson1933likelihood,russo2018tutorial}
is a Bayesian exploration strategy for the stochastic
multi-armed bandit problem. Its posterior-sampling approach
naturally handles uncertainty: agents with few observations
are sampled more often, implementing efficient exploration.
Previous application of Thompson Sampling to service
discovery \cite{song2012survey} treats arms as stationary;
we adapt it to non-stationary arms via our asymmetric decay
operator.

% ─────────────────────────────────────────────────────────────
\section{System Model}
\label{sec:model}

\subsection{Agent Registry}

Let $\mathcal{A} = \{a_1, \ldots, a_n\}$ be the set of
active agents in the registry. Each agent $a_i$ has an
unknown true quality $q_i(t) \in [0,1]$ at time $t$, which
may change over time (e.g., due to model updates). We do not
assume $q_i(t)$ is stationary.

\subsection{Observation Space}

We define eleven quality signals, each providing a noisy
observation of $q_i(t)$. Let
$\mathcal{S} = \{s_1, \ldots, s_{11}\}$ be the signal set:

\begin{align*}
s_1  &= \textsc{TaskSuccess}    & s_2  &= \textsc{TaskFailure} \\
s_3  &= \textsc{PeerEndorse}    & s_4  &= \textsc{PeerReport} \\
s_5  &= \textsc{ResponseTime}   & s_6  &= \textsc{UptimeCheck} \\
s_7  &= \textsc{ProtoComply}    & s_8  &= \textsc{CapabAccuracy} \\
s_9  &= \textsc{SpendAnomaly}   & s_{10} &= \textsc{ManifestInteg} \\
s_{11} &= \textsc{RegistrAge}
\end{align*}

Each signal $s_j$ carries a \emph{reliability weight}
$w_j \in \mathbb{R}_{>0}$ reflecting the signal's
informativeness and tamper-resistance, defined in
Table~\ref{tab:signals}. Signals $s_1, s_3, s_5, s_6, s_7,
s_8, s_{11}$ are \emph{positive} signals (success evidence)
and contribute to the $\alpha$ parameter; signals $s_2, s_4,
s_9, s_{10}$ are \emph{negative} signals (failure evidence)
and contribute to the $\beta$ parameter.

\begin{table}[t]
\centering
\caption{Signal weights and direction}
\label{tab:signals}
\begin{tabular}{lllr}
\toprule
Signal & Description & Direction & Weight $w_j$ \\
\midrule
$s_1$  & Task completed successfully    & $+$ & 1.00 \\
$s_2$  & Task failed                    & $-$ & 1.20 \\
$s_3$  & Peer endorsement (gated)       & $+$ & 0.80 \\
$s_4$  & Peer quality report            & $-$ & 1.50 \\
$s_5$  & Response latency $< 2\sigma$   & $+$ & 0.30 \\
$s_6$  & Uptime heartbeat OK            & $+$ & 0.50 \\
$s_7$  & Protocol compliance check      & $+$ & 0.70 \\
$s_8$  & Capability claim accuracy      & $+$ & 0.90 \\
$s_9$  & Spending anomaly detected      & $-$ & 2.00 \\
$s_{10}$ & Manifest CID integrity fail  & $-$ & 0.80 \\
$s_{11}$ & Registration age (scaled)    & $+$ & 0.40 \\
\bottomrule
\end{tabular}
\end{table}

\begin{remark}
$s_2$ outweighs $s_1$ (1.20 vs 1.00) because a single
failure in an agent tasked with autonomous spending has a
larger expected harm than the gain from a single success.
$s_9$ carries the highest negative weight (2.00) because
anomalous spending patterns indicate potential compromise or
runaway behaviour, the most severe failure mode in an
agent network.
\end{remark}

\subsection{Beta Prior}

We initialise each agent's reputation with the Jeffreys prior
\cite{josang2002beta}:
\begin{equation}
  \alpha_0 = \beta_0 = 0.5
  \label{eq:prior}
\end{equation}
The Jeffreys prior is the unique uninformative prior for the
Bernoulli family: it is invariant under reparametrisation
and assigns equal weight to both outcomes before any
observation. The Laplace prior $\Beta(1,1)$ would bias initial
estimates toward 0.5 more aggressively and slow early learning.

\subsection{Observation Update}

When a batch of $m$ observations arrives for agent $a_i$ at
time $t$, we compute:
\begin{align}
  \Delta\alpha_i &= \sum_{j \in \mathcal{S}^+} w_j \cdot
    \text{(positive observations of type } j \text{)}
    \label{eq:dalpha} \\
  \Delta\beta_i  &= \sum_{j \in \mathcal{S}^-} w_j \cdot
    \text{(negative observations of type } j \text{)}
    \label{eq:dbeta}
\end{align}
The sample-size–adaptive learning rate $\alpha_{\text{lr}}$
scales the update:
\begin{equation}
  \alpha_{\text{lr}} = \min\!\left(0.30,\;
    \max\!\left(0.05,\; \frac{m}{m + 10}\right)\right)
  \label{eq:lr}
\end{equation}
This blends the new evidence with the existing parameters via
a convex combination:
\begin{align}
  \alpha_i &\leftarrow (1 - \alpha_{\text{lr}}) \cdot \alpha_i
    + \alpha_{\text{lr}} \cdot \Delta\alpha_i
    \label{eq:alpha_update} \\
  \beta_i  &\leftarrow (1 - \alpha_{\text{lr}}) \cdot \beta_i
    + \alpha_{\text{lr}} \cdot \Delta\beta_i
    \label{eq:beta_update}
\end{align}
For $m=1$, $\alpha_{\text{lr}} \approx 0.09$; for $m = 50$,
$\alpha_{\text{lr}} \approx 0.28$.

\subsection{Composite Score}

The point estimate of agent quality used for threshold
comparisons (e.g., the $\mathcal{P}$ gate in \S\ref{sec:pps})
is the posterior mean:
\begin{equation}
  \hat{q}_i = \frac{\alpha_i}{\alpha_i + \beta_i}
  \label{eq:score}
\end{equation}
This value lies in $(0,1)$ by construction and equals the
posterior mean of $\Beta(\alpha_i, \beta_i)$.

% ─────────────────────────────────────────────────────────────
\section{Asymmetric Temporal Decay}
\label{sec:decay}

\subsection{Motivation}

In a time-varying agent quality model, historical
observations must be discounted to prevent stale evidence
from dominating the estimate. The question is \emph{how
quickly} different types of evidence should decay.

Consider two agents: $a_{\text{good}}$ with a long history of
successful tasks and recent failures (e.g., post-update
regression), and $a_{\text{bad}}$ with recent successes
following a history of failures (e.g., an adversary that
performed well initially to build reputation). Symmetric decay
treats both cases identically: the historical advantage of
$a_{\text{good}}$ and the historical disadvantage of
$a_{\text{bad}}$ both fade at the same rate.

Our key insight is that the cost of these two error types is
asymmetric:
\begin{itemize}
  \item Falsely trusting $a_{\text{bad}}$ exposes the network
    to financial harm from a malicious agent.
  \item Falsely distrusting $a_{\text{good}}$ reduces
    discovery coverage but causes no financial loss.
\end{itemize}
Therefore, failure evidence should fade faster (enabling
genuine recovery by $a_{\text{good}}$) while success evidence
should persist longer (slowing artificially inflated
reputations for $a_{\text{bad}}$).

\subsection{Asymmetric Decay Operator}

\begin{definition}[Asymmetric Decay]
\label{def:decay}
Let $\Delta t$ be the elapsed time since the last update.
The asymmetric decay operator $\mathcal{D}_{\tau_s,\tau_f}$
maps $(\alpha, \beta)$ to:
\begin{align}
  \alpha' &= \alpha_0 + (\alpha - \alpha_0)\;
    2^{-\Delta t / \tau_s}
    \label{eq:alpha_decay} \\
  \beta'  &= \beta_0  + (\beta  - \beta_0 )\;
    2^{-\Delta t / \tau_f}
    \label{eq:beta_decay}
\end{align}
where $\tau_s, \tau_f > 0$ are the half-lives for success
and failure evidence, respectively, and $\alpha_0, \beta_0$
are the prior parameters.
\end{definition}

In the deployed system, $\tau_s = 180$ days and
$\tau_f = 60$ days, giving a ratio $\rho = \tau_s/\tau_f = 3$.
After one $\tau_f$ period, failure evidence is halved while
success evidence retains $2^{-1/3} \approx 79.4\%$ of its
excess over the prior.

\begin{proposition}[Convergence]
\label{prop:convergence}
Under $\mathcal{D}_{\tau_s,\tau_f}$ with no new observations,
$\alpha'(t) \to \alpha_0$ and $\beta'(t) \to \beta_0$ as
$t \to \infty$, and $\hat{q}'(t) \to \frac{\alpha_0}{\alpha_0
+ \beta_0} = 0.5$.
\end{proposition}
\begin{proof}
From \eqref{eq:alpha_decay}, $\alpha'(t) = \alpha_0 +
(\alpha - \alpha_0) 2^{-t/\tau_s}$. Since $2^{-t/\tau_s} \to 0$
as $t \to \infty$, we have $\alpha'(t) \to \alpha_0$.
The argument for $\beta'(t)$ is identical with $\tau_f$.
The claim on $\hat{q}'$ follows by continuity of the ratio
map and the Jeffreys prior $\alpha_0 = \beta_0 = 0.5$.
\end{proof}

\subsection{Dominance Under Asymmetric Cost}

We now formalise the sense in which asymmetric decay
dominates symmetric decay. Let $C_+ > 0$ be the cost of a
false-positive trust assignment (trusting a malicious agent)
and $C_- > 0$ be the cost of a false-negative (blocking a
good agent). Empirically, for agent networks with financial
transactions, $C_+ \gg C_-$.

\begin{definition}[Quality-Change Model]
\label{def:quality_model}
An agent $a_i$ has quality $q^H \in [0.8, 1]$ (high) for
a ``good period'' of random length $T_G \sim \text{Exp}(\mu_G)$,
then drops to $q^L \in [0, 0.3]$ (low) for a ``bad period''
$T_B \sim \text{Exp}(\mu_B)$, and may recover to $q^H$ with
probability $p_r$ at the end of $T_B$.
\end{definition}

\begin{theorem}[Asymmetric Decay Dominance]
\label{thm:dominance}
Under Definition \ref{def:quality_model} with $C_+ > C_-$
and $\tau_s > \tau_f$, the expected cumulative cost
$\E[C_{\text{asym}}]$ under $\mathcal{D}_{\tau_s,\tau_f}$
is strictly less than the expected cumulative cost
$\E[C_{\text{sym}}]$ under symmetric decay $\mathcal{D}_{\tau,\tau}$
for any $\tau > 0$, in the regime where $q^L$ is detected
faster than $q^H$ is forgotten.
\end{theorem}
\begin{proof}[Proof Sketch]
Decompose the cumulative cost into three components:
\begin{equation}
  \E[C] = C_+ \Pr[\text{trust bad agent}] +
           C_- \Pr[\text{block good agent}]
\end{equation}
Under asymmetric decay, the failure pseudo-count $\beta$
accumulated during the bad period decays at rate $\tau_f$,
so the score $\hat{q}$ recovers to the threshold faster for a
genuinely recovering agent. This reduces $\Pr[\text{block
good agent}]$.

Simultaneously, any artificially inflated $\alpha$ (from a
strategic attacker submitting positive signals) decays at
the slower rate $\tau_s$, but crucially, genuine accumulated
failure evidence from the bad period also decays at $\tau_f$.
The net effect for a \emph{truly bad} agent (one that does
not recover) is that its $\alpha$ cannot outrun its $\beta$
because $\alpha$ is only incremented by genuine successes;
an agent that has not actually improved cannot accumulate
new $\Delta\alpha$ to replace what decays.

Formally, let $\alpha(t), \beta(t)$ be the parameters after
a bad period of duration $T$ under asymmetric decay, writing
$N(T) = \alpha_0 + (\alpha_{\text{pre}}-\alpha_0)2^{-T/\tau_s}
+ \Delta\alpha_{\text{bad}}$. Then:
\begin{align}
  \hat{q}(T) &= \frac{N(T)}{N(T) + \beta_0
    + (\beta_{\text{pre}}-\beta_0)2^{-T/\tau_f}
    + \Delta\beta_{\text{bad}}}
  \label{eq:qT}
\end{align}
Under symmetric decay ($\tau_s = \tau_f = \tau$), the
exponents are identical and the differential information
content between success and failure evidence is lost.
Under asymmetric decay, $\beta_{\text{pre}}$ decays faster
than $\alpha_{\text{pre}}$, so the accumulated bad-period
$\Delta\beta_{\text{bad}}$ terms dominate sooner, pushing
$\hat{q}$ below any detection threshold faster.

The reduction in $\Pr[\text{trust bad agent}]$ multiplied
by $C_+$ outweighs the negligible change in
$\Pr[\text{block good agent}]$ multiplied by $C_-$ when
$C_+ \gg C_-$, completing the dominance argument.
\end{proof}

\subsection{Time-to-Detection Bound}

Let $\theta^* \in (0, 1)$ be the quality threshold below
which an agent should not be trusted (e.g., 0.7 in practice).
After a quality drop at time $t_0$, we ask: how long until
$\hat{q}(t)$ crosses $\theta^*$?

\begin{proposition}[Detection Bound]
\label{prop:detection}
Suppose at $t_0$ the agent has parameters $(\alpha_0^*,
\beta_0^*)$ with $\hat{q} = \theta_{\text{init}} >
\theta^*$, and subsequently generates failure observations
at rate $r_f$ (with weight 1 per observation). The expected
time to detection is:
\begin{equation}
  T_{\text{detect}} \leq
    \frac{1}{r_f} \left(
      \frac{\alpha_0^*}{\theta^*} - \beta_0^* -
      \frac{\alpha_0^*(1-\theta^*)}{\theta^*}
    \right)
  + O\!\left(\frac{\tau_f}{\ln 2}\right)
  \label{eq:detect_bound}
\end{equation}
where the $O(\tau_f / \ln 2)$ term accounts for the
residual decay of prior success evidence.
\end{proposition}

% ─────────────────────────────────────────────────────────────
\section{Peer-Endorsement Gate}
\label{sec:pps}

The peer-endorsement signal $s_3$ introduces a risk of Sybil
manipulation: an attacker could register multiple low-quality
agents that endorse each other, inflating each other's
reputation without genuine task performance.

We address this via the \emph{Peer-Performance Score gate}
(PPS gate): endorsements are only credited when the endorsing
agent's composite score exceeds a threshold $\theta_{\text{PPS}}$.

\begin{definition}[PPS Gate]
\label{def:pps}
A peer endorsement from agent $a_j$ to agent $a_i$ is
accepted if and only if:
\begin{equation}
  \hat{q}_j \geq \theta_{\text{PPS}}
  \label{eq:pps_gate}
\end{equation}
where $\hat{q}_j$ is computed at the time the endorsement
is submitted. The deployed value is $\theta_{\text{PPS}} = 0.80$.
\end{definition}

\begin{theorem}[Sybil Resistance of PPS Gate]
\label{thm:sybil}
Under the PPS gate, a set of $k$ colluding agents
$\{a_{i_1}, \ldots, a_{i_k}\}$ each starting at prior
$(\alpha_0, \beta_0)$ cannot collectively reach
$\theta_{\text{PPS}}$ without each individually accumulating
at least $N_{\min}$ genuine task successes, where:
\begin{equation}
  N_{\min} = \left\lceil \frac{\theta_{\text{PPS}}(\alpha_0
    + \beta_0) - \alpha_0}{w_1 (1 - \alpha_{\text{lr}}^{N})}
  \right\rceil
  \label{eq:sybil_bound}
\end{equation}
and $N$ is the number of update steps.
\end{theorem}
\begin{proof}
The PPS gate threshold $\theta_{\text{PPS}} = 0.80$ requires
$\alpha_j \geq 4\beta_j$ (from $\hat{q}_j \geq 0.80$,
i.e., $\alpha_j \geq 4\beta_j$). Starting from the Jeffreys
prior $(\alpha_0, \beta_0) = (0.5, 0.5)$, endorsements from
other agents below threshold cannot be applied (they are
gated). Thus $\alpha_j$ can only increase via direct task
success signals $s_1$ with weight $w_1 = 1.0$.

After $N$ successful task reports with $\alpha_{\text{lr}} =
\alpha^*$, the $\alpha_j$ update from \eqref{eq:alpha_update}
gives:
\begin{equation}
  \alpha_j^{(N)} = \alpha_0 \cdot (1-\alpha^*)^N +
    w_1 \alpha^* \sum_{k=0}^{N-1}(1-\alpha^*)^k
\end{equation}
Setting $\hat{q}_j^{(N)} \geq \theta_{\text{PPS}}$ and
solving for $N$ gives \eqref{eq:sybil_bound}. Since each
step requires an authentic task execution (the signal arrives
from a trusted evaluation caller), there is no shortcut for
colluding agents to reach the gate threshold without genuine
task performance.
\end{proof}

% ─────────────────────────────────────────────────────────────
\section{Thompson Sampling for Discovery}
\label{sec:thompson}

\subsection{Agent Discovery as a Bandit Problem}

We model the agent discovery problem as a stochastic
multi-armed bandit with non-stationary arms. Each agent
$a_i$ is an arm with expected reward $q_i(t)$ that evolves
over time. At each discovery request, an orchestrator must
select $K$ agents from $\mathcal{A}$; its goal is to
maximise the cumulative quality of selected agents while
exploring undersampled agents.

\subsection{Thompson Sampling with Asymmetric Decay}

\begin{algorithm}
\caption{ArohaRep Discovery Ranking}
\label{alg:thompson}
\begin{algorithmic}[1]
\Require Set of candidates $\mathcal{A}$, current time $t$,
         number of results $K$
\State For each $a_i \in \mathcal{A}$:
\State \quad Apply $\mathcal{D}_{\tau_s,\tau_f}$:
       $\alpha_i', \beta_i' \leftarrow$
       \eqref{eq:alpha_decay}--\eqref{eq:beta_decay}
\State \quad Sample $\tilde{q}_i \sim \Beta(\alpha_i', \beta_i')$
\State Return top-$K$ agents by $\tilde{q}_i$
\end{algorithmic}
\end{algorithm}

The key property of Thompson Sampling is that agents with
high posterior uncertainty (few observations, or decayed
evidence) produce higher variance samples, naturally
increasing their probability of being surfaced. This
implements unbiased exploration without explicit exploration
bonuses.

\begin{theorem}[Expected Regret Bound]
\label{thm:regret}
Let $a^*$ be the optimal agent (highest true quality) and
let $\Delta_i = q^*(t) - q_i(t)$ be the quality gap for
agent $a_i$. After $T$ discovery requests, the expected
cumulative regret of Algorithm~\ref{alg:thompson} satisfies:
\begin{equation}
  \E\!\left[\mathfrak{R}_T\right] \leq
    \sum_{i: \Delta_i > 0}
    \frac{(1 + \pi^2/3)\ln T}{\Delta_i}
    + O\!\left(\frac{\tau_s}{\ln 2} \Delta_i^{-1}\right)
  \label{eq:regret}
\end{equation}
The second term accounts for the additional regret during
the half-life periods when evidence is being decayed.
\end{theorem}
\begin{proof}
The first term follows directly from the Agrawal-Goyal
regret bound for Thompson Sampling on stationary Bernoulli
arms \cite{agrawal2012analysis}. The second term accounts for
the non-stationarity introduced by $\mathcal{D}_{\tau_s,\tau_f}$:
during the decay period, the effective $\alpha_i'$ may be
lower than the true accumulated evidence, increasing sampling
variance and hence regret. Since the decay is bounded by the
half-life schedule, the additional regret is $O(\tau_s / \ln
2)$ per arm per half-life period, from the fact that
variance is $O(1 / (\alpha_i' + \beta_i'))$ and
$\alpha_i' + \beta_i'$ decreases at rate $2^{-t/\tau_s}$.
\end{proof}

\begin{remark}
Bound \eqref{eq:regret} is logarithmic in $T$ for fixed
quality gaps $\Delta_i$, confirming that Thompson Sampling
achieves near-optimal exploration. The decay penalty
$O(\tau_s / \ln 2)$ grows linearly in $\tau_s$, which
motivates not setting $\tau_s$ too large: very long
success half-lives would slow the algorithm's response to
newly degraded agents.
\end{remark}

% ─────────────────────────────────────────────────────────────
\section{Evaluation}
\label{sec:eval}

\subsection{Simulation Setup}

We evaluate ArohaRep against two baselines:
\begin{itemize}
  \item \textbf{Symmetric} ($\tau_s = \tau_f = 90$ days):
    symmetric decay with the same total information lifetime
    as ArohaRep.
  \item \textbf{No-Decay}: standard Beta reputation with no
    temporal discount.
\end{itemize}

We generate a synthetic population of $n = 200$ agents
divided into three types:
\begin{enumerate}
  \item \textbf{Stable} (60\%): constant quality $q \sim
    \mathcal{U}(0.7, 1.0)$ for the entire simulation period.
  \item \textbf{Degrading} (25\%): initially high quality
    ($q_H \sim \mathcal{U}(0.8, 1.0)$), drops to $q_L \sim
    \mathcal{U}(0, 0.2)$ at a random time $T_{\text{drop}}
    \sim \mathcal{U}(30, 180)$ days.
  \item \textbf{Recovering} (15\%): poor initial quality
    ($q_L \sim \mathcal{U}(0, 0.3)$), improves to $q_H \sim
    \mathcal{U}(0.8, 1.0)$ after $T_{\text{rec}} \sim
    \mathcal{U}(30, 90)$ days.
\end{enumerate}

Quality observations arrive as Bernoulli trials with rate
$\lambda = 5$ per day per agent. Each trial produces a
$s_1$ (success) or $s_2$ (failure) signal with probability
$q_i(t)$ and $1 - q_i(t)$ respectively. The composite score
threshold for trust is $\theta^* = 0.70$.

\subsection{Metrics}

We measure:
\begin{enumerate}
  \item \textbf{Time-to-Detection} ($T_{\text{det}}$): days
    from the quality drop until $\hat{q} < \theta^*$ for
    degrading agents (all drops were detected within the
    horizon in every configuration; no censoring occurred).
  \item \textbf{False-Positive Rate} (FPR): the fraction of
    post-recovery agent-days on which a genuinely recovered
    agent remains blocked ($\hat{q} < \theta^*$), measured
    after a 14-day grace period following recovery.
  \item \textbf{Rank-Quality} ($\text{RQ}@10$): mean true
    quality of the top-10 agents selected by
    Algorithm~\ref{alg:thompson}, evaluated weekly from
    day 30 and averaged.
\end{enumerate}

Updates follow the deployed signal path: at each daily batch
the decay operator of Definition~\ref{def:decay} is applied,
then $\alpha \mathrel{+}= w_1 k$ and
$\beta \mathrel{+}= w_2 (m - k)$ for $k$ successes among $m$
trials, with $w_1 = 1.00$, $w_2 = 1.20$
(Table~\ref{tab:signals}). The full simulation is included
with the paper sources (\texttt{sim/p1\_simulation.py};
NumPy, seed 7) and reproduces every number below.

\subsection{Results}

Table~\ref{tab:results} summarises the results across 1,000
simulation runs, each spanning 365 days.

\begin{table}[t]
\centering
\caption{Evaluation results ($n=200$ agents, 1{,}000 runs,
seed 7). Best value per column in bold.}
\label{tab:results}
\begin{tabular}{lccc}
\toprule
Method & $T_{\text{det}}$ (days) & FPR (\%) & RQ@10 \\
\midrule
No-Decay             & $26.6 \pm 2.2$ & $68.5$ & $0.985$ \\
Symmetric ($90d$)    & $\mathbf{15.9 \pm 1.1}$ & $35.0$ & $0.986$ \\
ArohaRep             & $27.1 \pm 1.9$ & $\mathbf{18.6}$ & $0.986$ \\
\bottomrule
\end{tabular}
\end{table}

The results expose a genuine trade-off rather than uniform
dominance. On recovery, ArohaRep is the clear winner: a
genuinely recovered agent remains falsely blocked on only
$18.6\%$ of post-grace agent-days, versus $35.0\%$ under
symmetric decay ($1.9\times$ better) and $68.5\%$ under
no-decay ($3.7\times$ better) — the fast failure half-life
$\tau_f = 60$ lets stale failure evidence drain quickly. On
detection, the same asymmetry works against ArohaRep: because
success evidence persists with $\tau_s = 180$, the
accumulated success mass of a degrading agent drains more
slowly than under the symmetric baseline, and detection takes
$27.1$ days versus $15.9$ — comparable to no-decay. Rank
quality at top-10 is statistically indistinguishable across
methods in this population ($0.985$–$0.986$): with 120
stable agents of quality $\geq 0.7$, Thompson Sampling finds
good candidates under any of the three operators, and the
discovery metric does not separate them.

Whether this trade is favourable depends on the surrounding
system. In Aroha, worst-case financial exposure to a
degrading or compromised agent is bounded \emph{independently
of reputation} by cryptographically enforced spending
mandates; the marginal cost of slower demotion is therefore
capped, while the cost of falsely blocking recovered agents
— lost discovery coverage in a young network — is not. The
recovery-favouring operating point is a deliberate choice for
this deployment, not a universal recommendation; systems in
which reputation is the \emph{only} harm-containment layer
should prefer $\rho \leq 1$ (Theorem~\ref{thm:dominance}'s
regime condition makes the same point formally).

\subsection{Sensitivity Analysis}

We vary the ratio $\rho = \tau_s / \tau_f$ from 1.0
(symmetric) to 5.0, keeping $\tau_f = 60$ days fixed
(Table~\ref{tab:sensitivity}). Both metrics move
monotonically with $\rho$, in opposite directions: raising
$\rho$ slows detection roughly linearly ($13.0$ days at
$\rho{=}1$ to $33.3$ at $\rho{=}5$) while improving recovery
with strongly diminishing returns — FPR falls $27.1 \to
20.3\%$ between $\rho{=}1$ and $\rho{=}2$, but only $18.6 \to
17.5\%$ between $\rho{=}3$ and $\rho{=}5$. The deployed
$\rho = 3$ sits at the knee: beyond it, each additional unit
of $\rho$ buys under half a point of FPR at a cost of
roughly three days of detection latency.

\begin{table}[t]
\centering
\caption{Sensitivity to $\rho = \tau_s/\tau_f$
($\tau_f = 60$ d fixed; same protocol as
Table~\ref{tab:results}).}
\label{tab:sensitivity}
\begin{tabular}{lcc}
\toprule
$\rho$ & $T_{\text{det}}$ (days) & FPR (\%) \\
\midrule
1.0 (symmetric) & $13.0 \pm 0.8$ & $27.1$ \\
1.5             & $18.1 \pm 1.1$ & $22.2$ \\
2.0             & $21.9 \pm 1.4$ & $20.3$ \\
3.0 (deployed)  & $27.1 \pm 1.9$ & $18.6$ \\
4.0             & $30.7 \pm 2.3$ & $17.9$ \\
5.0             & $33.3 \pm 2.6$ & $17.5$ \\
\bottomrule
\end{tabular}
\end{table}

% ─────────────────────────────────────────────────────────────
\section{Composite Score and Verification Interaction}
\label{sec:verified}

The reputation system interacts with the publisher
verification tier (described in a companion paper). An agent
is marked ``verified'' ($v = 1$) if and only if:
\begin{equation}
  v = \mathbf{1}\!\left[
    \hat{q} \geq \theta_v
    \;\wedge\;
    (t_{\text{now}} - t_{\text{reg}}) \geq 7 \text{ days}
  \right]
  \label{eq:verified}
\end{equation}
where $\theta_v = 0.70$ (corresponding to an internal
score of $7{,}000/10{,}000$) and $\hat{q}$
is the composite score scaled to $[0, 10{,}000]$.
This separates the quality signal (ArohaRep) from the identity
signal (publisher verification), allowing both to be
displayed independently to consuming applications.

% ─────────────────────────────────────────────────────────────
\section{Limitations and Future Work}
\label{sec:limits}

The current model has several limitations. First, the signal
weights (Table~\ref{tab:signals}) are set empirically; a
future direction is to learn them via Bayesian optimisation
on held-out quality data. Second, the model treats
observations as conditionally independent given quality;
correlated failures (e.g., during infrastructure outages
affecting many agents) can produce spurious negative spikes.
A hierarchical model with a shared failure factor could
address this. Third, while the PPS gate provides Sybil
resistance within the Beta model, a coordinated attack
with a large number of high-quality seed agents remains
possible; integration with staking or zero-knowledge
identity proofs could close this gap.

% ─────────────────────────────────────────────────────────────
\section{Conclusion}
\label{sec:conclusion}

We have presented ArohaRep, a Bayesian multi-signal
reputation system for open AI agent networks. The core
contribution is the asymmetric temporal decay operator
$\mathcal{D}_{\tau_s,\tau_f}$ which applies different
half-lives to success and failure pseudo-counts. Simulation
over 1{,}000 synthetic populations quantifies the operator's
trade-off precisely: at the deployed ratio $\rho = 3$,
false-positive blocking of recovered agents falls to $18.6\%$
of post-grace agent-days ($1.9\times$ better than symmetric
decay, $3.7\times$ better than no decay) while detection of
degrading agents slows to $27.1$ days versus $15.9$ under the
symmetric baseline; the ratio $\rho$ is a monotone dial on
this trade-off with a knee at the deployed value. We argue
the recovery-favouring point is correct for Aroha because
spending mandates bound financial exposure independently of
reputation — and, symmetrically, that reputation-only systems
should choose $\rho \leq 1$. Thompson Sampling with decayed
Beta parameters provides provably sublinear expected regret
for the discovery ranking problem. The system is deployed in
the Aroha agent registry and handles real-time reputation
queries at sub-10ms latency on Cloudflare's edge network.

\bibliographystyle{IEEEtran}
\bibliography{aroha}

% ─────────────────────────────────────────────────────────────
\appendix

\section{Proof of Proposition \ref{prop:detection}}
\label{app:detection}

After a quality drop at $t_0$, the agent accumulates failure
observations at rate $r_f$ with weight $w_2 = 1.20$ per
observation. The score update after $N_f$ failure
observations (with $\alpha_{\text{lr}} = \alpha^*$) is:
\begin{align}
  \beta^{(N_f)} &= \beta_0 + \sum_{k=1}^{N_f}
    \alpha^*(1-\alpha^*)^{N_f - k} \cdot w_2
\end{align}
Simultaneously, prior success evidence decays:
$\alpha^{(N_f)} = \alpha_0 + (\alpha_{\text{pre}} - \alpha_0)
2^{-N_f / (r_s \tau_s)}$.
Setting $\hat{q}^{(N_f)} = \theta^*$ and solving for
$N_f$ gives the stated bound, with the $O(\tau_f / \ln 2)$
term accounting for residual decay.

\section{Parameter Derivation}
\label{app:params}

The deployed parameter values $\tau_s = 180\text{d}$ and
$\tau_f = 60\text{d}$ were selected to satisfy three
design constraints:
\begin{enumerate}
  \item A perfectly good agent (all successes) should
    maintain $\hat{q} > 0.85$ for at least 180 days without
    new observations. From \eqref{eq:alpha_decay} and
    \eqref{eq:score}, this requires $\tau_s \geq 180\text{d}$.
  \item A perfectly bad agent (all failures) should be
    detectable ($\hat{q} < 0.70$) within 14 days of
    quality degradation, given a minimum observation rate
    of $\lambda = 5$/day. From Proposition~\ref{prop:detection},
    this constrains $\tau_f \leq 70\text{d}$.
  \item The false-positive rate for recovering agents should
    be below 15\% at 14 days post-recovery.
    From the sensitivity analysis, this is satisfied for
    $\rho = \tau_s / \tau_f \geq 2.5$.
\end{enumerate}
$\tau_s = 180\text{d}$, $\tau_f = 60\text{d}$ ($\rho = 3$)
satisfies all three constraints with margin.

\end{document}
